\chapter{Integration with the Design}
\label{ch:integration}

The preceding chapters justify each tool on its own merits. This chapter shows
that the set of them fits the design already agreed, by mapping every component
of that design to the tool that implements it, and every requirement to the tools
that serve it.

References to diagrams are to the UML Design Document, and requirement
identifiers follow the scheme used there: \req{SA} for state machine to activity
diagram, \req{AC} for activity diagram to code, \req{CS} for code to state
diagram, \req{FS} for stretch functional requirements, and \req{NFR} for
non-functional requirements of the minimum viable product.

\section{Components to Tools}

Each component of the block diagram is implemented by a named tool in a named
project. Nothing in the design is left without an implementer, and no tool
appears that the design does not call for.

\begin{longtable}{L{0.21\textwidth}L{0.24\textwidth}L{0.45\textwidth}}
	\caption{Block diagram components and what implements them}\label{tab:components}\\
	\toprule
	\textbf{Component} & \textbf{Tool} & \textbf{Project and notes} \\
	\midrule
	\endfirsthead
	\toprule
	\textbf{Component} & \textbf{Tool} & \textbf{Project and notes} \\
	\midrule
	\endhead

	Repository Client & GENESYS Client API & \code{Pentateuch.Genesys}. The only
	component that references the vendor assemblies, which is what the block
	diagram asserts and what makes the licence constraint local rather than
	global. \\
	\addlinespace

	Model Mapper & None beyond the language & \code{Pentateuch.Genesys}. Converts
	between Client API entity types and the shared model. Deliberately no library:
	the mapping is specific to this schema and a general-purpose mapper would
	obscure it. \\
	\addlinespace

	Code Extractor & Roslyn & \code{Pentateuch.Translator}. Uses
	\code{Microsoft.CodeAnalysis.Workspaces.MSBuild} to open a solution and
	\code{Microsoft.CodeAnalysis.CSharp} for the semantic model and control flow
	analysis. \\
	\addlinespace

	Abstraction & None beyond the language & \code{Pentateuch.Translator}.
	Operates on the shared model rather than on either end's types, which is what
	allows one component to serve both directions as the block diagram shows. \\
	\addlinespace

	Code Generator & Roslyn & \code{Pentateuch.Translator}. Builds syntax trees
	and formats them through
	\code{Microsoft.CodeAnalysis.CSharp.Workspaces}. \\
	\addlinespace

	GENESYS Model & GENESYS & External. Not built by this project. \\
	\addlinespace

	Code Repository & Git & External. The tool reads and writes files; version
	control operations remain the user's. \\
	\bottomrule
\end{longtable}

The class diagram divides the same system differently, by representation rather
than by component, and the tools line up with it just as cleanly. The
\code{Absyn} class at the centre of that diagram is \code{Pentateuch.Absyn}, the
.NET Standard 2.0 library. The \code{C\#} class is Roslyn's view of source. The
\code{Genesys Project} class with its \code{Function}, \code{Item} and
\code{State} entities is the Client API's view of the repository. That the two
decompositions resolve to the same three projects is a reasonable check that
neither is arbitrary.

\section{Requirements to Tools}

\begin{longtable}{L{0.14\textwidth}L{0.33\textwidth}L{0.43\textwidth}}
	\caption{Requirements and the tools that serve them}\label{tab:req-tools}\\
	\toprule
	\textbf{Req} & \textbf{Tool} & \textbf{How} \\
	\midrule
	\endfirsthead
	\toprule
	\textbf{Req} & \textbf{Tool} & \textbf{How} \\
	\midrule
	\endhead

	\req{SA-1} to \req{SA-3} & GENESYS Client API only & Both ends are model
	entities, so this translation never leaves \code{Pentateuch.Genesys} and needs
	no compiler at all. It is the one direction Roslyn plays no part in. \\
	\addlinespace

	\req{AC-1} & Client API, then Roslyn & The Client API reads the activity, the
	shared model carries it, Roslyn emits the source. \\
	\addlinespace

	\req{AC-2} & Roslyn & Generating a class per activity node with function
	prototypes attached is syntax tree construction, which is what Roslyn's
	factories are for. \\
	\addlinespace

	\req{AC-3} & Roslyn & Forks, loops and merges become language constructs,
	emitted as syntax nodes rather than assembled as text. \\
	\addlinespace

	\req{AC-4} & Roslyn & Compiling the generated source answers whether it
	builds. This is the reason the \emph{Verify Generated Code Compiles} action
	exists in the export activity diagram, and it is available only because the
	generator is a compiler rather than a text templater. \\
	\addlinespace

	\req{CS-1} & Roslyn & \code{MSBuildWorkspace} loads the repository and
	\code{Microsoft.CodeAnalysis.CSharp} produces the control flow graph. \\
	\addlinespace

	\req{CS-2} & Roslyn & Following calls to a designated depth requires knowing
	which method a call resolves to, which is the semantic model rather than the
	syntax tree. \\
	\addlinespace

	\req{CS-3} & Roslyn & Branch coverage is read off the control flow graph the
	compiler already builds, rather than being reconstructed. \\
	\addlinespace

	\req{CS-4} & Roslyn & Recursive and mutually recursive calls are detectable
	because the semantic model identifies the target method, so the queue in the
	import activity diagram can refuse to enqueue one twice. \\
	\addlinespace

	\req{CS-5} & Roslyn & Reporting unused variables, methods and classes needs
	reference resolution across the whole solution, which is
	\code{SymbolFinder}. \\
	\addlinespace

	\req{CS-6} & Roslyn & Deriving state names from the code uses identifiers and
	structure that the syntax tree preserves. \\
	\addlinespace

	\req{FS-1} & GENESYS scripting & Not adopted. The committed deliverable runs
	beside GENESYS rather than inside it, and the architecture is arranged so that
	a script could later start and control it. \\
	\addlinespace

	\req{FS-3} & Roslyn and Git & Rename reconciliation needs both: Roslyn to
	recognise that a symbol is the same one under a new name, and Git's rename
	detection for files that moved. \\
	\addlinespace

	\req{NFR-2} & The runtime split & That a failure in the tool cannot terminate
	GENESYS follows from the two programs being separate processes, which is the
	decision in Chapter~\ref{ch:runtime} rather than a property of any one
	tool. \\
	\addlinespace

	\req{NFR-4} & Pinned package versions & Determinism requires that every
	developer runs the same Roslyn, since analysis behaviour changes between
	releases. \\
	\addlinespace

	\req{NFR-5} & GENESYS & Windows-only is inherited from the host, not
	chosen. \\
	\bottomrule
\end{longtable}

\section{How Design Decisions Constrained the Tools}

Four decisions were made before the toolchain was settled, and each closed off
options rather than merely informing them. Recording which came first matters,
because it is the difference between a design that follows from its tools and
tools that follow from a design.

\paragraph{The process boundary came first.}
The block diagram separates the components that run with GENESYS from those that
do not. That boundary is why there are three projects rather than one, why one of
them targets .NET Standard 2.0, and why a current Roslyn is available at all. Had
the design placed everything in one process, the tool list would have been
shorter and considerably worse.

\paragraph{Confining the vendor dependency came first.}
The block diagram puts the Client API behind a single component. That is what
makes \req{NFR-2} achievable and what allows most of the codebase to build
without a licence, which in turn is why a limited number of activations is an
inconvenience rather than a blocker.

\paragraph{The abstract syntax as pivot came first.}
The class diagram places \code{Absyn} between the code representation and the
model representation. \code{Pentateuch.Absyn} is that class made real, and
.NET Standard 2.0 was selected because it is the only target both halves can
reference. The library exists to satisfy the diagram, not the other way round.

\paragraph{Generated views ruled out a whole category of tool.}
Diagrams in GENESYS are generated from entities and relationships rather than
stored as pictures. The tool therefore creates entities and never draws anything,
which means no graph layout or diagramming library appears anywhere in this
document. That absence is a consequence of a design fact, not an oversight
